\documentclass[10pt,a4paper]{amsart}
\usepackage[margin=19mm]{geometry}
\usepackage{amsmath,amssymb,mathtools}
\usepackage[scr=rsfs,cal=euler]{mathalfa}
\usepackage{xfrac}
\usepackage[unicode,hidelinks,bookmarks=false]{hyperref}
\newcommand{\ie}{i.e.\ }
\newcommand{\cf}{cf.\ }
\newcommand{\resp}{respectively\ }
\newcommand{\Subsection}{\subsection}
\providecommand{\nobreakdash}{\nobreak\hbox{-}}
\setlength{\emergencystretch}{2em}
\multlinegap0pt
\usepackage{bm}
\usepackage{bbm}
\usepackage{dsfont}
\usepackage{xspace}
\usepackage{cancel}
\usepackage{mathtools}
\usepackage{amsthm}
\makeatletter
\@ifundefined{theorem}{\newtheorem{theorem}{Theorem}}{}
\@ifundefined{proposition}{\newtheorem{proposition}[theorem]{Proposition}}{}
\makeatother
\DeclareMathOperator{\var}{Var}
\newcommand{\E}{\mathbb{E}}
\title[On absolute continuity of inhomogeneous and contracting on a]{On absolute continuity of inhomogeneous and contracting on average self-similar measures}
\author{Samuel Kittle, Constantin Kogler}
\date{Source: arXiv:2409.18936v4 (CC BY 4.0)}
\begin{document}
\maketitle

\noindent\emph{Source and licence.} On absolute continuity of inhomogeneous and contracting on average self-similar measures. Version arXiv:2409.18936v4 (CC BY 4.0), distributed under the Creative Commons Attribution 4.0 International licence, as recorded in the arXiv licence metadata for this exact version and shown on the abstract page https://arxiv.org/abs/2409.18936v4. Full rights notice: licenses/arxiv-2409.18936-CC-BY-4.0.txt.\par\smallskip

\noindent\textit{Editorial scope.} Counted unit: the Proposition whose source label is \texttt{VarSumAdds} in the article (source0.tex lines 1476--1481, source label \texttt{VarSumAdds}), delivered together with the article's own complete original proof of it (source0.tex lines 1483--1543, one \texttt{proof} environment of 61 lines). No mathematics is written, repaired or re-derived: statement and proof are the article's own text line for line. Required context retained uncounted: none. Every definition, display and label of those lines is retained unchanged. Omitted from this package and not redistributed: 1--1475, 1482--1482, 1544--2459. Nothing inside a retained excerpt is omitted or altered: every retained body is delivered exactly as the article's lines are written, so no omission receipt of any kind accompanies this package. Citations: the retained excerpts carry no citation command at all, so no bibliography entry of the article is transcribed and no reference is redistributed. Reference adaptation: none was needed and none was made; every reference inside the delivered unit resolves inside it, so no unresolved reference or citation is produced. Printed slips kept literal: no printed word, symbol, hypothesis or number of the counted statement or of its proof was altered. Layout and class: the article's own class is \texttt{amsart} with packages including inputenc, babel, hyperref, color, enumerate, fancyhdr, amsmath, amscd, amsfonts, amsthm, amstext, amssymb, amsbsy, mathrsfs; the wrapper is the lane's pinned \texttt{amsart} editorial wrapper at 10pt on A4, which restates the theorem-like environments the retained text needs and adds the article's own macro definitions that the retained text needs (\texttt{\textbackslash E}, \texttt{\textbackslash var}), with extra wrapper packages (bm, bbm, dsfont, xspace, cancel, mathtools). The delivered numbering is the unit's own. Assets: no retained span contains a figure environment, a graphics-inclusion command, a PDF-page inclusion, a table or a TikZ picture; no image file is copied into this package, the package declares no asset dependency and no image file is redistributed with it. Licence: version arXiv:2409.18936v4 of this article is distributed under the Creative Commons Attribution 4.0 International licence, as recorded in the arXiv licence metadata for this exact version and shown on the abstract page https://arxiv.org/abs/2409.18936v4. A full rights notice is delivered at licenses/arxiv-2409.18936-CC-BY-4.0.txt. Characters: every retained excerpt is pure ASCII, so the delivered engine, XeLaTeX with its default Latin Modern text font, reports no missing character.
\begin{proposition} \label{VarSumAdds}
    Let $\mu$ be a probability measure on $G$. Let $R > 1$ be such that $\rho(g) \in [R^{-1}, R]$ for every $g \in \mathrm{supp}(\mu)$. Let $n_1, n_2, K \in \mathbb{Z}_{\geq 0}$ with $n_2, K > 0$ and let $\kappa_1, \kappa_2, r \in (0, 1)$. Let $A > 0$ and let $M \geq R$. Then
    \begin{align*}
        \MoveEqLeft V(\mu,  n_1 + n_2, K, R^{-1} M^{-1} \kappa_1 \kappa_2, A; r) \\&\geq V(\mu, n_1, K, \kappa_1, A; r) + V(\mu,  n_2, K, \kappa_2, A; M \kappa_1^{-1} r).
    \end{align*} 
\end{proposition}

\begin{proof}
    For $j \in \{ 1,2 \}$ let $\gamma_1^{(j)}, \gamma_2^{(j)}, \dots$ be a sequence of i.i.d.\ samples from $\mu$ defined on the probability space $\left(\Omega_{(j)}, \mathscr{F}_{(j)}, \mathbb{P}_{(j)}\right)$. Let $\hat{\gamma}_1, \hat{\gamma}_2, \dots$ be a sequence of i.i.d.\ samples from $\mu$ defined on the probability space $\left(\hat{\Omega}, \hat{\mathscr{F}}, \hat{\mathbb{P}}\right)$. Consider the product probability space $$\left(\Omega, \mathscr{F}, \mathbb{P}\right) = \left( \Omega_1 \times \hat{\Omega} \times \Omega_2, \mathscr{F}_1 \times \hat{\mathscr{F}} \times \mathscr{F}_2, \mathbb{P}_1 \times \hat{\mathbb{P}} \times \mathbb{P}_2 \right).$$ 

    Let $\left(\gamma_i^{(1)}, S_i^{(1)}, T_i^{(1)}, f_i^{(1)}, U_i^{(1)}, h_i^{(1)}, \mathscr{A}_i^{(1)}, m_i^{(1)}\right)$ be a proper decomposition for $(\mu, k_1, K, \kappa_1, A)$ at scale $r$ defined on the probability space $\left(\Omega^{(1)}, \mathscr{F}^{(1)}, \mathbb{P}^{(1)}\right)$ such that $\sum_{i = 1}^{k_1} m_i^{(1)}$ approaches $V(\mu,n_1,K,\kappa_1,A;r)$ and  $$ \rho(f_1^{(1)}h_1^{(1)}\cdots f_{k_1}^{(1)}h_{k_1}^{(1)}) \geq \kappa_1.$$ 

    Given $\omega_1 \in \Omega_1$ and $\hat{\omega} \in \hat{\Omega}$, let $\tau = \tau(\omega_1, \hat{\omega})$ be given by $$\tau = \min \{k \in \mathbb{Z}_{\geq 0}: \rho( f^{(1)}_1 h^{(1)}_1 f^{(1)}_2 h^{(1)}_2 \dots f^{(1)}_{k_1} h^{(1)}_{k_1} \hat{\gamma}_1 \hat{\gamma}_2 \dots \hat{\gamma}_k) < M^{-1} \kappa_1  \} $$ and let $\hat{\rho} = \rho( f^{(1)}_1 h^{(1)}_1 f^{(1)}_2 h^{(1)}_2 \dots f^{(1)}_{k_1} h^{(1)}_{k_1} \hat{\gamma}_1 \hat{\gamma}_2 \dots \hat{\gamma}_\tau)$ such that $$\hat{\rho} \in [M^{-1} R^{-1} \kappa_1 , M^{-1} \kappa_1 ]. $$

    Now given $\omega_1 \in \Omega_1$ and $\hat{\omega} \in \hat{\Omega}$, let $\left(\gamma_i^{(2)}, S_i^{(2)}, T_i^{(2)}, f_i^{(2)}, U_i^{(2)}, h_i^{(2)}, \mathscr{A}_i^{(2)}, m_i^{(2)}\right)$ be a proper decomposition for $(\mu, k_2, K, \kappa_2, A)$ at scale $M\kappa_1^{-1}r$ defined on the probability space $\left(\Omega^{(2)}, \mathscr{F}^{(2)}, \mathbb{P}^{(2)}\right)$ such that $\sum_{i = 1}^{k_2} m_i^{(2)}$ approaches $V(\mu,n_2,K,\kappa_2,A;M\kappa_1^{-1}r)$ and  $$\rho(f_1^{(1)}h_1^{(1)}\cdots f_{k_2}^{(1)}h_{k_2}^{(1)}) \geq \kappa_2.$$ 

    We now concatenate the two decompositions as follows. Let $\gamma_1, \gamma_2, \dots$ be the sequence of random variables on the probability space $\left(\Omega, \mathscr{F}, \mathbb{P}\right)$ defined by $$\gamma_i = \begin{cases}
        \gamma_i^{(1)} & \text{if } i \leq T^{(1)}_{k_1} \\
        \hat{\gamma}_{i - T^{(1)}_{k_1}} & \text{if } i > T^{(1)}_{k_1} \text{ and } i \leq T^{(1)}_{k_1} + \tau \\
        \gamma_{i - T^{(1)}_{k_1} - \tau} & \text{if } i > T^{(1)}_{k_1} + \tau.
    \end{cases}$$
    Clearly these are i.i.d.\ samples from $\mu$. For $i = 1, 2, \dots, k_1 + k_2$ we define $S_i$ by
    \begin{equation*}
        S_i =
        \begin{cases}
            S_i^{(1)} & \text{if } i \leq k_1\\
            S_{i-k_1}^{(2)} + T^{(1)}_{k_1} + \tau & \text{if } i > k_1
        \end{cases}
    \end{equation*}
    and we define $T_i$ analogously. We define $f_i$ by
    \begin{equation*}
        f_i =
        \begin{cases}
            f_i^{(1)} & \text{if } i \leq k_1\\
            \hat{\gamma}_1 \dots \hat{\gamma}_{\tau} f_{1}^{(2)} & \text{if } i = k_1 + 1 \\
            f_{i-k_1}^{(2)} & \text{if } i > k_1 + 1.
        \end{cases}
    \end{equation*}
    We define $U_i$ by
    \begin{equation*}
        U_i =
        \begin{cases}
            U_i^{(1)} & \text{if } i \leq k_1\\
            U_{i-k_1}^{(2)} & \text{if } i > k_1.
        \end{cases}
    \end{equation*}
    and define $h_i$ and $m_i$ analogously. Finally we define $\mathscr{A}_i$ by
    \begin{equation*}
        \mathscr{A}_i =
        \begin{cases}
            \mathscr{A}_i^{(1)} \times \hat{\Omega} \times \Omega^{(2)} & \text{if } i \leq k_1\\
            \mathscr{A}_{k_1}^{(1)} \times \hat{\mathscr{F}} \times \mathscr{A}_{i-k_1}^{(2)} & \text{if } i > k_1.
        \end{cases}
    \end{equation*}
    It is easy to check that $\left(\gamma_i, S_i, T_i, f_i, U_i, h_i, \mathscr{A}_i, m_i\right)$ is a proper decomposition for $(\mu, R, k_1 + k_2, K, R^{-1} M^{-1} \kappa_1 \kappa_2, A)$ at scale $r$ and it holds that $$\sum_{i=1}^{k_1+k_2} m_i = \sum_{i=1}^{k_1} m_i^{(1)} + \sum_{i=1}^{k_2} m_i^{(2)}.$$ Indeed, we note that for $i > k_2$ we have that since $M\kappa_1^{-1} \leq \hat{\rho}^{-1}$,
    \begin{align*}
        |U_i| = |U_{i-k_1}^{(2)}| &\leq \rho(f_1^{(2)}h_1^{(2)}f_2^{(2)}h_2^{(2)}\cdots h_{i-k_1 -1}^{(2)} f_{i-k_1}^{(2)})^{-1} M\kappa_1^{-1}r \\
        &\leq \hat{\rho}^{-1}\rho(f_1^{(2)}h_1^{(2)}f_2^{(2)}h_2^{(2)}\cdots h_{i-k_1 -1}^{(2)} f_{i-k_1}^{(2)})^{-1} r \\
        &= \rho(f_1h_1f_2h_2\cdots h_{i-1} f_{i})^{-1}r.
    \end{align*}
    Similarly, for $i > k_2 + 1$ and using that $\hat{\rho}^2M^2\kappa_1^{-2} \leq 1$,
    \begin{align*}
        \MoveEqLeft \E\left[   \frac{\var(\rho(f_i) U(f_i) U_i b(h_i)|\mathscr{A}_i)}{\rho(f_1h_1 f_2h_2 \cdots f_{i-1}h_{i-1})^{-2} r^2} \,|\, \mathscr{A}_{i-1}  \right] \\
        &= \E\left[   \frac{\var(\rho(f^{(2)}_{i -k_1}) U(f^{(2)}_{i -k_1}) U^{(2)}_{i-k_1} b(h^{(2)}_{i - k_1})|\mathscr{A}_i)}{\hat{\rho}^{-2}\rho(f_1^{(2)}h_1^{(2)}f_2^{(2)}h_2^{(2)}\cdots h_{i-k_1}^{(2)})^{-2} r^2} \,|\, \mathscr{A}_{i-1}  \right] \\
        &\geq  \E\left[   \frac{\var(\rho(f^{(2)}_{i -k_1}) U(f^{(2)}_{i -k_1}) U^{(2)}_{i-k_1} b(h^{(2)}_{i - k_1})|\mathscr{A}_i)}{\hat{\rho}^{-2}\rho(f_1^{(2)}h_1^{(2)}f_2^{(2)}h_2^{(2)}\cdots h_{i-k_1}^{(2)})^{-2} \hat{\rho}^2 M^2 \kappa_1^{-2} r^2} \,|\, \mathscr{A}_{i-1}  \right] \\
        &\geq m_{i - k_1}^{(2)}I.
    \end{align*} The remainder of the properties are straightforward to check. 
\end{proof}

\end{document}
